\chapter{Decisioni di progetto e razionale}
\label{cap:sdd-decisioni}

Le decisioni che seguono hanno determinato la forma del sistema. Per ciascuna
sono riportati il problema, la scelta, le alternative considerate e le
conseguenze accettate. Sono raccolte qui perché una decisione motivata è
riesaminabile, mentre una decisione implicita si può solo subire.

\section{DP\_1 — La sessione è un token firmato in un cookie httpOnly}

\textbf{Problema.} Riconoscere l'utente fra una richiesta e l'altra
proteggendo il meccanismo di riconoscimento.

\textbf{Decisione.} Un JWT firmato dal server, con scadenza breve, depositato
in un cookie \id{httpOnly}, \id{SameSite=Lax} e \id{Secure} in produzione.

\textbf{Alternative.} Il token in un'intestazione \id{Authorization},
conservato dal client, sarebbe più comodo da provare con strumenti da riga di
comando, ma richiede che il JavaScript di pagina possa leggerlo: una
vulnerabilità di scripting basterebbe a rubarlo. Le sessioni sul server
renderebbero la revoca immediata, al costo di uno stato condiviso e di un
archivio di sessioni, contro \id{DG\_6}.

\textbf{Conseguenze.} Il server resta privo di stato. In cambio, un token non
può essere revocato prima della scadenza: è il motivo per cui la scadenza è
breve. Il cookie viaggia cross-origin, quindi il CORS deve dichiarare
un'origine esplicita e ammettere le credenziali.

\section{DP\_2 — La correzione del quiz avviene solo sul server}

\textbf{Problema.} Impedire che l'esito di una verifica possa essere
falsificato.

\textbf{Decisione.} La rappresentazione del quiz inviata al client non
contiene l'indicazione della risposta corretta. Il client raccoglie le scelte
e le consegna; il server corregge e restituisce esito e soluzioni.

\textbf{Alternative.} Correggere nel browser eviterebbe un giro di rete e
darebbe un riscontro istantaneo, ma richiederebbe di trasmettere le
soluzioni: chiunque potrebbe leggerle prima di rispondere.

\textbf{Conseguenze.} L'invariante è garantito dalla struttura del dato, non
dalla disciplina di chi scrive il codice: la funzione di serializzazione non
espone il campo, quindi non c'è modo di dimenticarsene. Il costo è una
richiesta HTTP al momento della consegna.

\section{DP\_3 — Le risposte sono associate per identificativo}

\textbf{Problema.} Correlare le scelte dell'utente alle domande del quiz.

\textbf{Decisione.} Il client invia coppie
\id{(domandaId, rispostaId)}; il server le indicizza per identificativo di
domanda.

\textbf{Alternative.} Un elenco ordinato di identificativi di risposta,
correlato per posizione all'ordine delle domande sul server, sarebbe più
compatto. Introdurrebbe però un accoppiamento implicito fra l'ordine di
presentazione nell'interfaccia e quello di memorizzazione: basterebbe
mostrare le domande in un ordine diverso — per esempio mescolandole a ogni
tentativo — perché la correzione risultasse errata senza alcun segnale.

\textbf{Conseguenze.} L'ordine di presentazione diventa una libera scelta
dell'interfaccia. Una risposta riferita a una domanda che non appartiene al
quiz viene semplicemente ignorata.

\section{DP\_4 — Il conteggio dei quiz superati è derivato}

\textbf{Problema.} Sapere quanti quiz un utente ha superato, dato che il
valore determina l'assegnazione dei badge.

\textbf{Decisione.} Il valore è ricalcolato contando gli svolgimenti con
esito positivo, e poi memorizzato nella colonna \file{utente.quiz\_superati}
come denormalizzazione di comodo.

\textbf{Alternative.} Incrementare un contatore a ogni superamento è più
economico, ma lega il valore alla storia delle operazioni anziché allo stato
del sistema: ripetere un quiz già superato lo incrementerebbe di nuovo, e i
badge finirebbero assegnati per traguardi mai raggiunti.

\textbf{Conseguenze.} Un'interrogazione di conteggio in più a ogni consegna,
in cambio di un valore che non può divergere dalla realtà. La colonna
memorizzata resta comoda per le letture e viene riscritta solo quando il
conteggio effettivo differisce.

\section{DP\_5 — La valutazione media è mantenuta dal database}

\textbf{Problema.} Mostrare nel catalogo la media dei feedback di ogni
tutorial senza calcolarla a ogni lettura.

\textbf{Decisione.} Tre trigger su \file{feedback} ricalcolano
\file{tutorial.valutazione} a ogni inserimento, modifica o cancellazione.

\textbf{Alternative.} Calcolarla nel servizio richiederebbe di ricordarsene
in tre punti diversi, e due scritture concorrenti potrebbero lasciare un
valore errato. Calcolarla a ogni lettura del catalogo costringerebbe a un
raggruppamento su tutta la tabella dei feedback.

\textbf{Conseguenze.} Il valore è coerente per costruzione, anche di fronte a
scritture concorrenti. In cambio, una parte del comportamento del sistema
vive nel database, non è visibile nel codice applicativo e non è coperta
dalla suite di test: è dichiarata qui e nel Test Document.

\section{DP\_6 — Quiz, domande e risposte formano un solo aggregato}

\textbf{Problema.} Stabilire quali entità abbiano un DAO proprio.

\textbf{Decisione.} Un unico \id{QuizDao} governa quiz, domande e risposte.
Non esistono DAO per domanda e risposta.

\textbf{Alternative.} Un DAO per tabella è la scelta più immediata, ma
introduce tre problemi. La lettura di un quiz richiederebbe
un'interrogazione per ogni domanda e una per ogni risposta. Le domande
potrebbero essere scritte senza il loro quiz, in stati che il dominio non
ammette. E i DAO di domanda e risposta sarebbero troppo sottili per essere
verificati in modo significativo da soli, pur restando il punto in cui è più
facile sbagliare l'ordine di un parametro.

\textbf{Conseguenze.} La lettura di un quiz è una sola interrogazione
(\id{RNF\_P\_1}). L'aggregato è sempre coerente: non esiste un modo di
scrivere una domanda senza il suo quiz.

\section{DP\_7 — Il client non condivide codice con il server}

\textbf{Problema.} Evitare che i due lati divergano sui tipi scambiati.

\textbf{Decisione.} Il client ridichiara in \file{src/types} le strutture che
riceve dall'API. Nessun import attraversa il confine.

\textbf{Alternative.} Importare le entità del server direttamente nel client
eviterebbe la duplicazione ed è tecnicamente possibile con una regola di
compilazione dedicata. Il client dipenderebbe però da classi che contengono
logica di dominio e accesso ai dati, e non potrebbe essere costruito senza il
server.

\textbf{Conseguenze.} I due progetti si compilano e si distribuiscono in modo
indipendente. Il prezzo è una duplicazione delle definizioni di trasporto,
accettata perché DTO ed entità hanno ragioni di esistere diverse e non devono
cambiare insieme.

\section{DP\_8 — L'assistente conversazionale è un componente esterno}

\textbf{Problema.} Offrire assistenza immediata a un pubblico inesperto.

\textbf{Decisione.} Il sistema può incorporare un widget di dialogo fornito
da terzi, abilitato solo se configurato. Non esiste logica conversazionale
propria.

\textbf{Alternative.} Realizzare un assistente interno richiederebbe
competenze, infrastruttura e dati di addestramento estranei agli obiettivi
del progetto; un assistente inaffidabile, su questo pubblico, sarebbe
dannoso più che utile.

\textbf{Conseguenze.} La funzionalità è dichiarata fuori perimetro nel RAD.
L'applicazione funziona identicamente senza il widget, e nessun dato
dell'utente gli viene trasmesso dal sistema.
